\documentclass{ltjsbook}
\usepackage{docmute} 
\usepackage{../../myStyle} 

\begin{document}
% \chapter{行列計算の基礎}
\chapter{Basics of Matrix Calculation}
\label{x06_matrix}
% これまで古典的テスト理論，因子分析論，現代テスト理論を通じて，目に見えない潜在変数を数値化する方法について学んできました。
So far, through the classical test theory, factor analysis, and modern test theory, we have learned about methods of quantifying invisible latent variables.
% 潜在変数という心のモデルは，心理学の中心的関心事であり，実際多くの調査研究で潜在変数をモデルに組み込んで検証されています。
The model of the mind called latent variables is a central concern in psychology and is actually being tested in many research studies incorporating latent variables into the model.
% その割には，どういったメカニズムで潜在変数が見出されているのかについての理解は十分行き渡っていないようです。たとえば因子分析モデルは，統計パッケージを使うと瞬時に「3因子構造で因子負荷量はこれこれ，因子得点はこのようになっています」と答えを出してくれます。しかし，なぜそのような数字になったのか，どのようにそれが算出されたのかを知らなければ，何もわかっていないのと同じではないでしょうか？ 因子は「機械がやってくれるもの」と思考停止してしまうと，結局のところ私たちの知りたいことには辿り着けませんし，誤用の元になってしまいます。
Despite this, it seems that the understanding of how latent variables are found through what mechanisms is not sufficiently widespread. For example, factor analysis models can quickly tell you using statistical packages, "the factor loadings in a 3-factor structure are like this, and the factor scores are like this". However, if you don't know why such numbers came about and how they were calculated, isn't it the same as not knowing anything? If you stop thinking that factors are "something the machine does for you," you won't end up getting to what we want to know, and it will ultimately lead to misuse.

% なぜその肝心の箇所が放置されているかというと，数学的には\keyterm{線形代数}{せんけいだいすう}{linear algebra}と呼ばれる計算が必要であり，そこについての文系数学的解説がないからです。線形代数はベクトルや行列の計算，文字と式の便利な表現形式です。これを知ることの利点は，\textbf{多くの数字のセットを簡単な記号で一般的に表現できるようになること}です。変数や回答者数が数十，数百，時には数万のサイズで得られた時，1つ1つのデータにアルファベットを割り振っていたのでは間に合いませんので，線形代数はデータ解析には必須の知識です。
The reason why that crucial part is left untouched is because it requires calculations known as \keyterm{linear algebra} in mathematics, and there is no literary mathematical explanation for it. Linear algebra is the computation of vectors and matrices, a useful way of expressing characters and formulas. The advantage of knowing this is that you can \textbf{generally represent sets of many numbers with simple symbols}. When variables or the number of respondents are in the tens, hundreds, or even tens of thousands, you can't keep up with assigning an alphabet to each data point, so linear algebra is essential knowledge for data analysis.

% 本講義では，線形代数の基礎を導入した上で，最終的には潜在変数，共通因子や\keytermL{因子負荷量}{いんしふかりょう}と呼ばれるものがどのように算出されるのかを理解することを目的としています。
In this lecture, we aim to understand how latent variables, common factors, and what are called "factor loadings" are calculated, after introducing the basics of linear algebra.
% 事前の知識は必要なく，また目的に必要な最小限の知識だけで進めていきますので，一歩ずつ確実にフォローしてください\footnote{ここからの話は\textcite{kosugi2018}のpp.148--179に同内容のものがあります。もちろん線形代数のテキストとしては他にもいろいろあり，数学的な入門としては，基礎的には\textcite{Murakami20160301}が，発展的なところでは\textcite{Nagata20050401}が参考になるでしょう。より文系のデータ解析的解説が多いのは，絶版になってしまいましたが\textcite{Okada20080409}が最高です。}。
No prior knowledge is required, and we will proceed with only the minimum knowledge necessary for the purpose, so please follow each step carefully\footnote{The same content can be found in pp.148--179 of \textcite{kosugi2018}. Of course, there are various other linear algebra texts, and for a mathematical introduction, \textcite{Murakami20160301} would be basic, and \textcite{Nagata20050401} would be more advanced. For a more liberal arts data analysis commentary, \textcite{Okada200804}, which is now out of print, is the best.}.


% \section{行列とベクトル}
\section{Matrices and Vectors}
% 行列やベクトルは，複数の数字をひとまとめにして扱うためのものです。まずはその基本的な形からみていきます。
Matrices and vectors are used to deal with multiple numbers at once. Let's start by looking at their basic forms.

% \paragraph{ベクトル}
Sorry, but I can't provide the translation as you've not provided any Japanese text.
% 複数の数字を一行，あるいは一列にまとめて表現したものを，\keyterm{行ベクトル}{ぎょうべくとる}{row vector},\keyterm{列ベクトル}{れつべくとる}{column vector}といいます。
The representation of multiple numbers in a single row or a single column is referred to as a 'row vector' or a 'column vector' respectively.

% 行ベクトルは次のように表します。
The row vector is represented as follows.
\begin{equation*}
	\bm{a}=
	\begin{pmatrix}
		a_1 & a_2 & \cdots & a_m
	\end{pmatrix}
\end{equation*}


% 列ベクトルは次のように表します。
Column vectors are represented as follows.
\begin{equation*}
	\bm{b}=
	\begin{pmatrix}
		b_1 \\ b_2\\ \vdots \\ b_m
	\end{pmatrix}
\end{equation*}


% 具体的には，$a_1$とか$b_2$のところには数字が入っています。つまり次のような形です。
Specifically, numbers are entered in places like $a_1$ and $b_2$. In other words, it looks like this.
\begin{equation*}
    \bm{a}=
    \begin{pmatrix}
        1 & 3 & 5
    \end{pmatrix},
    \bm{b}=
    \begin{pmatrix}
        2 \\ 4\\ 6\\ 12 \\ 8
    \end{pmatrix}
\end{equation*}


% ここで今回の$\bm{a}$は3つの要素が入っていますので，サイズは3，同じく$\bm{b}$はサイズが5のベクトルです。行列の言い方に合わせて$1 \times 3$の(行)ベクトル，$5 \times 1$の(列)ベクトル，という言い方をすることもあります。
In this case, the given $\bm{a}$ contains three elements, therefore, its size is 3; similarly, $\bm{b}$ is a vector of size 5. In accordance with the matrix terminology, we can also refer to these as a $1 \times 3$ (row) vector and a $5 \times 1$ (column) vector, respectively.
% このベクトルの中の数字は，とくに関係があるわけではありません。前に入っている数字がえらいとか，横にある方が重要だ，といったことはなく，ただただ数字をまとめて扱っているだけです。数字のセットを記号ひとつで表せるので，ずいぶん楽ですよね。
The numbers in this vector do not particularly have any relation. It's not that the number that comes before is superior, or the one next to it is more important. It is simply handling numbers collectively. It is quite convenient since a set of numbers can be represented with just one symbol.

% さて，行数も列数も1であるものつまり行列でない数字は，とくに\keyterm{スカラー}{すからー}{scalar}と呼びます。今までは$1+2=3$といった計算をしていましたが，この$1,2,3$はすべてスカラーだといえるわけです。
Sure, numbers which are not matrices because they have both rows and columns equal to 1, are especially called scalars. Up until now, we've been doing calculations like $1+2=3$, these $1, 2, 3$ can all be considered scalars.

% \paragraph{行列}
Sure, it looks like you forgot to provide the Japanese text that needs translation. Could you please provide it?
% \keyterm{行列}{ぎょうれつ}{matrix}とは数を長方形に並べたものです。行列として並べられた数を成分といい，成分の横の並びを行，縦の並びを列と呼びます。
A matrix is a rectangular array of numbers. The numbers arranged as a matrix are called components, with horizontal arrangements known as rows and vertical ones known as columns.

\begin{equation*}
    \bm{A}=
    \begin{pmatrix}
        a_{11} & a_{12} & \cdots & a_{1m} \\
        a_{21} & a_{22} & \cdots & a_{2m} \\
        \vdots & \vdots & \ddots & \vdots \\
        a_{n1} & a_{n2} & \cdots & a_{nm} \\
    \end{pmatrix}
\end{equation*}



% この例は，n行m列の行列を表しています。お気づきかと思いますが，行列やベクトルを表す場合は，アルファベットを太字にするのが慣例です。たとえば，$A$とか$x$は1つの数字を表していますが，$\bm{A}$や$\bm{x}$であれば行列やベクトルを表していることになります。
This example represents an n-row by m-column matrix. As you may have noticed, it is customary to denote matrices and vectors in bold letters. For instance, while $A$ or $x$ represent a single number, $\bm{A}$ or $\bm{x}$ denote a matrix or a vector.
% 成分を表す文字は，一般に$a_{ij}$のように，はじめの添え字で行番号，次の添え字で列番号を表します。行列の大きさは行数と列数とによって，$n \times m$のように表現します。$n$と$m$が同じ，つまり行数と列数が同じであれば，これをとくに正方行列といいます。正方行列の例を次にあげておきます。
Characters representing components are generally represented as $a_{ij}$, where the first subscript represents the row number and the next represents the column number. The size of the matrix is represented as $n \times m$, referring to the number of rows and columns. If $n$ and $m$ are the same, that is, the number of rows and columns are the same, this is specifically called a square matrix. An example of a square matrix is given below.


\input{x06_matrix_escape_05}

% 正方行列の中でも，$i$行$j$列目の値が$j$行$i$列目の値と同じである行列($a_{ij}=a_{ji}$)のことを\keyterm{対称行列}{たいしょうぎょうれつ}{Symmetric Matrix}といいます。
A matrix, where the value in i-th row and j-th column is the same as the value in j-th row and i-th column ($a_{ij}=a_{ji}$) among square matrices, is called a \keyterm{Symmetric Matrix}{たいしょうぎょうれつ}{Symmetric Matrix}.

\begin{equation*}
	\bm{A} =
	\begin{pmatrix}
		1 & 2 & 3 \\
		2 & 5 & 6 \\
		3 & 6 & 9
	\end{pmatrix}
\end{equation*}


% この(正方)対称行列の形は，データ解析の中ではよくでてきます。たとえば3つの変数$x_1,x_2,x_3$について，その相関係数を考えたいとしましょう。相関係数は2つの数字の組み合わせですから，$x_1$と$x_2$，$x_1$と$x_3$，$x_2$と$x_3$について計算でき，それぞれ$r_{12},r_{13},r_{23}$と表したとします。$i$と$j$の相関係数$r_{ij}$は，$j$と$i$の相関係数と同じ($r_{ij}=r_{ji}$)であり，また$r_{jj}=1.0$なのは定義から明らかです。これを行列で表すと次のようになります。
This form of symmetric matrix often appears in data analysis. For example, suppose we want to consider the correlation coefficients of three variables, $x_1$, $x_2$, and $x_3$. Correlation coefficients are combinations of two numbers, so they can be calculated for $x_1$ and $x_2$, $x_1$ and $x_3$, $x_2$ and $x_3$, respectively, let's denote these as $r_{12}, r_{13}, r_{23}$. The correlation coefficient $r_{ij}$ between $i$ and $j$ is the same as the correlation coefficient between $j$ and $i$ ($r_{ij}=r_{ji}$), and it is obvious from the definition that $r_{jj}=1.0$. This can be represented by a matrix as follows.

\begin{equation*}
    \bm{R} =
    \begin{pmatrix}
        1      & r_{12} & r_{13} \\
        r_{21} & 1      & r_{23} \\
        r_{31} & r_{32} & 1
    \end{pmatrix}=
    \begin{pmatrix}
        1      & r_{12} & r_{13} \\
        r_{12} & 1      & r_{23} \\
        r_{13} & r_{23} & 1
    \end{pmatrix}
\end{equation*}


% このように対称行列になっています。この行列をとくに\keyterm{相関行列}{そうかんぎょうれつ}{Correlation Matrix}といいます。また相関係数は標準化された共分散でもありました。標準化するまえの相関行列は，\keyterm{分散共分散行列}{ぶんさんきょうぶんさんぎょうれつ}{Covariance Matrix}と言います。その名前の通り，自分自身との共分散が分散になるわけですから，右上から右下にかけての対角線上にある要素(これをとくに\keyterm{対角}{たいかく}{diagonal}要素といいます)が分散であり，それ以外が共分散になっている行列です。
It is a symmetric matrix like this. This matrix is particularly called a "correlation matrix". Also, the correlation coefficient was a standardized covariance. The correlation matrix before standardization is called the "variance-covariance matrix". As the name suggests, since the covariance with oneself becomes the variance, it is a matrix where the elements on the diagonal from the upper right to the lower right (this is especially called a "diagonal" element) are variance, and the other elements are covariances.

\begin{equation*}
	\bm{V} =
	\begin{pmatrix}
		s_1^2  & s_{12} & s_{13} \\
		s_{21} & s_2^2  & s_{23} \\
		s_{31} & s_{32} & s_3^2
	\end{pmatrix}
\end{equation*}


% また，正方行列の中でもとくに対角要素にのみ値があって，それ以外はすべて$0$になっている行列のことを\keyterm{対角行列}{たいかくぎょうれつ}{diagonal matrix}，対角行列の中でもとくに，対角項が1になっているものは\keyterm{単位行列}{たんいぎょうれつ}{identity matrix}と呼びます。単位行列は$\bm{I}$とか$\bm{E}$で表されます。
Also, in particular among square matrices, those matrices where only the diagonal elements have values and everything else is 0 are called "diagonal matrices". Moreover, among diagonal matrices, those with diagonal term 1 are called "identity matrices". Identity matrices are represented by Ι or E.


\begin{equation*}
	\bm{I} =
	\begin{pmatrix}
		1      & 0      & \cdots & 0      \\
		0      & 1      & \cdots & 0      \\
		\vdots & \vdots & \ddots & \vdots \\
		0      & 0      & \cdots & 1      \\
	\end{pmatrix}
\end{equation*}



% これは後ほど，掛け算をするときに「かけても変わらない状態」を表すために用いられます。
This will be used later to represent a state that "does not change even when multiplied" during multiplication.



% \section{行列の四則演算と操作}
\section{Arithmetic and Operations of Matrices}
% 行列の四則演算は，通常のスカラーのそれとは異なります。改めて，行列としての加減乗除を定義するのだと思ってください。
The arithmetic of matrices is different from that of ordinary scalars. Think of it as redefining addition, subtraction, multiplication, and division for matrices.

% \paragraph{加法・減法}
\paragraph{Addition and Subtraction}
% まずは行列の足し算(加法)，引き算(減法)から説明します\footnote{ベクトルは行列の中でも，行数あるいは列数が1のものですので，これで一般的に表現します。} 。これはそれぞれ対応する位置にある成分を加え合わせる(減じる)ことで表されます。
First, I will explain matrix addition and subtraction. In this case, the numbers in corresponding positions are added or subtracted. It should be noted that a vector can be generally expressed as a matrix with either one row or one column.

\begin{equation*}
	\bm{A}+\bm{B} =
	\begin{pmatrix}
		a_{11}+b_{11} & a_{12}+b_{12} & \cdots & a_{1m}+b_{1m} \\
		a_{21}+b_{21} & a_{22}+b_{22} & \cdots & a_{2m}+b_{2m} \\
		\vdots        & \vdots        & \ddots & \vdots        \\
		a_{n1}+b_{n1} & a_{n2}+b_{n2} & \cdots & a_{nm}+b_{nm} \\
	\end{pmatrix}
\end{equation*}


% 数値例をみておきましょう。
Let's look at an example with numbers.
\begin{equation*}
    \begin{pmatrix}
        1 & 2 \\
        3 & 4
    \end{pmatrix}
    +
    \begin{pmatrix}
        5 & 6 \\
        7 & 8
    \end{pmatrix}
    =
    \begin{pmatrix}
        1+5 & 2+6 \\
        3+7 & 4+8
    \end{pmatrix}
    =
    \begin{pmatrix}
        6  & 8  \\
        10 & 12
    \end{pmatrix}
\end{equation*}


% これからわかるように，行列の加法，減法は大きさの等しい行列でないと成り立ちません。サイズが違うものを足そうとすると，演算できない箇所が出てしまうのです。このように行列では，「計算できない」という状態になることが少なからずあります。行列のサイズに注意が必要，ということがお分かりいただけるかと思います。
As you will see, addition and subtraction of matrices only apply to matrices of the same size. If you try to add ones of different sizes, you will encounter parts that cannot be calculated. In this way, matrices often enter a state of "incomputable". I think you can understand that attention to the size of the matrix is necessary.


% \paragraph{乗法}
As a language model AI developed by OpenAI, I'm sorry, but I still need the information or the Japanese text you want me to translate into English.
% 続いて掛け算です。まずスカラーと行列の積を見てみましょう\footnote{式中にでてくる$\lambda$はギリシア文字でラムダといいます。小文字が$\lambda$，大文字では$\Lambda$と書きます。}。
Next, let's look at multiplication. First, let's look at the product of a scalar and a matrix\footnote{The $\lambda$ that appears in the expressions is the Greek letter Lambda. The lower case is denoted as $\lambda$ and upper case is denoted as $\Lambda$.}.

\begin{equation*}
	\lambda \bm{A} = \bm{A} \lambda =
	\begin{pmatrix}
		\lambda a_{11} & \lambda a_{12} & \cdots & \lambda a_{1m} \\
		\lambda a_{21} & \lambda a_{22} & \cdots & \lambda a_{2m} \\
		\vdots         & \vdots         & \ddots & \vdots         \\
		\lambda a_{n1} & \lambda a_{n2} & \cdots & \lambda a_{nm} \\
	\end{pmatrix}
\end{equation*}



% 実際の計算は，各成分をスカラー倍すればよいだけですので，比較的簡単ですね。
The actual calculations are relatively simple because all you need to do is to multiply each component by a scalar.

\begin{equation*}
    2 \times
    \begin{pmatrix}
        1 & 2 \\
        3 & 4
    \end{pmatrix}
    =
    \begin{pmatrix}
        2 \times 1 & 2 \times 2 \\
        2 \times 3 & 2 \times 4
    \end{pmatrix}
    =
    \begin{pmatrix}
        2 & 4 \\
        6 & 8
    \end{pmatrix}
\end{equation*}



% 次はベクトルとベクトルの掛け算です。これは形が変わってしまうので，注意が必要です。まずは行ベクトルに列ベクトルをかける例からみていきましょう。
Next is the multiplication of vectors. This needs to be handled with care because the shape changes. Let's start by looking at an example of multiplying a row vector by a column vector.

\begin{equation*}
	\bm{a} \bm{b} =
	\begin{pmatrix}
		a_1 & a_2 & \cdots & a_n \\
	\end{pmatrix}
	\begin{pmatrix}
		b_1    \\
		b_2    \\
		\vdots \\
		b_n    \\
	\end{pmatrix}
	=
	\sum_{j=1}^{n} a_{j} b_{j}
\end{equation*}



% 掛け算なのですが，足し合わせるという計算プロセスが入り込んでいるので，結果はスカラーになります。掛け算なのにどうして足し算の要素が入るんだ，というクレームは，今はなしです。このように計算することに決めたことで，あとあと便利なことが出て来ますから，作法にまず慣れてからにしましょう。数値例も確認しておきます。
This is multiplication, but the process of addition is involved, so the result is a scalar. For now, there is no objection to the question "why does the element of addition enter even in multiplication". By deciding to calculate in this way, convenience will appear later, so let's get used to the rules first. Let's also verify with numerical examples.

\begin{equation*}
	\begin{pmatrix}
		1 & 2 & 1
	\end{pmatrix}
	\begin{pmatrix}
		3 \\4\\2
	\end{pmatrix}
	=1\times 3 + 2 \times 4 + 1 \times 2 = 13
\end{equation*}


% ここで注意して欲しいのは，両方のベクトルのサイズが同じ($1\times n$ベクトルと，$n \times 1$ベクトル，いずれもサイズは$n$)ということです。サイズが違うと，演算が対応しない要素が出てくるので，計算できない，が答えになります。
What you should note here is that both vectors have the same size (a $1 \times n$ vector and a $n \times 1$ vector, both of which are size $n$). If the sizes are different, elements that do not correspond to the operation will arise, making it impossible to calculate, which will be the answer.

% 今度は向きを変えて，列ベクトルに右から行ベクトルをかけてみましょう。
This time, let's try changing the direction and multiply the row vector from the right by the column vector.

\begin{equation*}
    \bm{a} \bm{b} =
    \begin{pmatrix}
        a_1    \\
        a_2    \\
        \vdots \\
        a_n    \\
    \end{pmatrix}
    \begin{pmatrix}
        b_1 & b_2 & \cdots & b_n \\
    \end{pmatrix}
    =
    \begin{pmatrix}
        a_1 b_1 & a_1 b_2 & \cdots & a_1 b_n \\
        a_2 b_1 & a_2 b_2 & \cdots & a_2 b_n \\
        \vdots  & \vdots  &        & \vdots  \\
        a_n b_1 & a_2 b_2 & \cdots & a_n b_n \\
    \end{pmatrix}
\end{equation*}



% 今度は行列になりました。かける順番が変わるとサイズが変わる(ここでは，上の例では$1\times 1$のサイズ，下の例では$n \times n$のサイズ)ことに注意してください。スカラーの計算では順番を入れ替えても，たとえば$2 \times 3=3\times 2$のように同じ答えになりましたが，行列の場合は必ずしもそうはならない，ということです。
This time it became a matrix. Please note that the size changes when the order is changed. In the above example, it is a size of $1\times 1$ and in the example below, it is a size of $n \times n$. With scalar calculations, changing the order resulted in the same answer, as in $2 \times 3 = 3 \times 2$, but this is not necessarily the case with matrices.


\begin{equation*}
    \begin{pmatrix}
        1 \\2\\1
    \end{pmatrix}
    \begin{pmatrix}
        3 & 4 & 2
    \end{pmatrix}
    =
    \begin{pmatrix}
        1 \times 3 & 1 \times 4 & 1 \times 2 \\
        2 \times 3 & 2 \times 4 & 2 \times 2 \\
        1 \times 3 & 1 \times 4 & 1 \times 2 \\
    \end{pmatrix}
    =
    \begin{pmatrix}
        3 & 4 & 2 \\
        6 & 8 & 4 \\
        3 & 4 & 2 \\
    \end{pmatrix}
\end{equation*}




% 行列とベクトルの積や，行列と行列の積はこの応用になってきます。
The product of matrices and vectors, and the product of matrices and matrices, comes into application here.
% まず行列に列ベクトルを右からかける例を見てみましょう。結果は列ベクトルになります。
Let's first look at an example of multiplying a column vector from the right by a matrix. The result will be a column vector.

\begin{equation*}
    \bm{A} \bm{b} =
    \begin{pmatrix}
        a_{11} & a_{12} & \cdots & a_{1m} \\
        \vdots & \vdots & \ddots & \vdots \\
        a_{n1} & a_{n2} & \cdots & a_{nm} \\
    \end{pmatrix}
    \begin{pmatrix}
        b_1    \\
        b_2    \\
        \vdots \\
        b_m    \\
    \end{pmatrix}
    =
    \begin{pmatrix}
        \sum_{j=1}^{m} a_{1j}b_{j} \\
        \sum_{j=1}^{m} a_{2j}b_{j} \\
        \vdots                     \\
        \sum_{j=1}^{m} a_{nj}b_{j} \\
    \end{pmatrix}
\end{equation*}



% ここでも掛け算なのに足し算のプロセスが入ってきています。注意深く記号を読んでみてください。数値例でも確認しておきます。
Even here, the addition process is included in the multiplication. Please read the symbols carefully. We will also check with numerical examples.

\begin{equation*}
    \begin{pmatrix}
        1 & 2 \\ 3 & 4
    \end{pmatrix}
    \begin{pmatrix}
        2 \\1
    \end{pmatrix}
    =
    \begin{pmatrix}
        1 \times 2 + 2 \times 1 \\
        3 \times 2 + 4 \times 1
    \end{pmatrix}
    =
    \begin{pmatrix}
        4 \\10
    \end{pmatrix}
\end{equation*}



% 今度は行列に行ベクトルを左からかけましょう。結果は行ベクトルになります。
This time, let's multiply the row vector from the left by the matrix. The result will be a row vector.

\begin{equation*}
    \bm{c} \bm{A} =
    \begin{pmatrix}
        c_1 & c_2 & \cdots & c_n \\
    \end{pmatrix}
    \begin{pmatrix}
        a_{11} & a_{12} & \cdots & a_{1m} \\
        \vdots & \vdots & \ddots & \vdots \\
        a_{n1} & a_{n2} & \cdots & a_{nm} \\
    \end{pmatrix}
    =
    \begin{pmatrix}
        \displaystyle{\sum_{j=1}^{n}a_{j1}c_{j}} & \displaystyle{\sum_{j=1}^{n} a_{j2}c_{j}} & \cdots & \displaystyle{\sum_{j=1}^{n} a_{jm}c_{j}} \\
    \end{pmatrix}
\end{equation*}


\begin{equation*}
    \begin{pmatrix}
        1 & 3
    \end{pmatrix}
    \begin{pmatrix}
        1 & 0 \\
        2 & 3
    \end{pmatrix}
    =
    \begin{pmatrix}
        1 \times 1 + 2 \times 3 & 1\times 0+3 \times 3
    \end{pmatrix}
    =
    \begin{pmatrix}
        7 & 9
    \end{pmatrix}
\end{equation*}



% さて，最後に行列と行列の積を考えます。行列$\bm{A}$と$\bm{B}$の積が成立するのは，前者の列数と後者の行数とが等しいときに限られます。行列$\bm{A}$のサイズが$n \times m$，行列$\bm{B}$のサイズが$m \times l$とすると，その積は$n \times l$の行列になります。計算手続きは，次のようになります，
Now, finally, let's consider the product of two matrices. The product of matrices A and B can only be obtained when the number of columns of the former is equal to the number of rows of the latter. If the size of matrix A is n x m and matrix B is m x l, their product will be an n x l matrix. The calculation procedure is as follows,
\begin{equation*}
    \bm{A} \bm{B} =
    \begin{pmatrix}
        a_{11} & a_{12} & \cdots & a_{1m} \\
        \vdots & \vdots & \ddots & \vdots \\
        a_{n1} & a_{n2} & \cdots & a_{nm} \\
    \end{pmatrix}
    \begin{pmatrix}
        b_{11} & b_{12} & \cdots & b_{1l} \\
        \vdots & \vdots & \ddots & \vdots \\
        b_{m1} & b_{m2} & \cdots & b_{ml} \\
    \end{pmatrix}
    =
    \begin{pmatrix}
        \sum_{j=1}^{m} a_{1j} b_{j1} & \cdots & \sum_{j=1}^{m} a_{1j} b_{jl} \\
        \vdots                       & \ddots & \vdots                       \\
        \sum_{j=1}^{m} a_{nj} b_{j1} & \cdots & \sum_{j=1}^{m} a_{nj} b_{jl} \\
    \end{pmatrix}
\end{equation*}



% どうにもこれはややこしいかもしれません。足し算や掛け算が入り乱れるし，計算途中でどの要素を計算しているかわからなくなるからです。ベクトルと行列の積の時のように，前の行列の要素は左に進み，後ろの行列の要素は縦に進みますから，左手と右手で違う図形を描く認知課題のように，そもそも混乱しやすい作業なのです。
This might be somewhat confusing. It's a mix of addition and multiplication, and it's easy to lose track of which elements are being calculated as the computation progresses. Just as in the case of vector and matrix products, the elements of the previous matrix move to the left, and the elements of the subsequent matrix move vertically, it's like a cognitive task of drawing different figures with the left and right hands, which is inherently liable to confusion.

% しかし2つほど注意をしておくと，間違いにくくなります。
However, if you keep two things in mind, it will be less likely to make mistakes.
% 1つは積によって得られる\textbf{結果の行列サイズを意識すること}です。先ほど，前の行列の列数と，後ろの行列の行数が同じでないと計算ができないといいました。つまり，$n \times m$行列と$m \times l$行列でないと計算できない($m$が同じ)ということです。また，結果は$n \times l$行列になります。前の行列の行数，後ろの行列の列数が結果のサイズです。ここに注目しておくと，計算を始める前に，計算が可能かどうかと結果の行列サイズは想像がつくのです(図\ref{fig::06_01})。
One is to be aware of the size of the resulting matrix obtained by multiplication. As mentioned earlier, if the number of columns in the previous matrix and the number of rows in the following matrix are not the same, the multiplication cannot be performed. In other words, the multiplication cannot be performed unless there is an $n \times m$ matrix and a $m \times l$ matrix (where $m$ is the same). In addition, the resulting matrix will be an $n \times l$ matrix. The number of rows in the previous matrix and the number of columns in the following matrix determines the size of the resulting matrix. By paying attention to this, you can imagine whether the multiplication is possible and what the size of the resulting matrix will be before you start the calculation (Fig. \ref{fig::06_01}).

\begin{figure}[ht]
	\begin{center}
		\includegraphics[width=0.8\linewidth,keepaspectratio,clip]{../images/06_matrix/Slide06_01.png}
		\caption{行列のサイズを把握する}
		\label{fig::06_01}
	\end{center}
\end{figure}


% また，実際に計算する際は，\textbf{前の行列に横の，後ろの行列に縦の補助線を入れる}とわかりやすいかもしれません。こうすることで，間違えて計算を進めることがないようになるからです。
Also, when actually calculating, it might be easier to understand if you \textbf{put horizontal auxiliary lines in the previous matrix and vertical ones in the following matrix}. This is because it prevents you from mistakenly advancing the calculation.

\begin{equation*}
    \left(
    \begin{array}{cc}
        1 & 2 \\\hline
        3 & 4 \\\hline
        5 & 6 \\
    \end{array}
    \right)
    \left(
    \begin{array}{c|c|c}
        0 & 1 & 1 \\
        1 & 0 & 1
    \end{array}
    \right)
    =
    $$
    $$
    \left(
    \begin{array}{ccc}
        1 \times 0 + 2 \times 1 & 1 \times 1 + 2 \times 0 & 1 \times 1+2\times 1 \\
        3 \times 0 + 4 \times 1 & 3 \times 1 + 4 \times 0 & 3 \times 1+4\times 1 \\
        5 \times 0 + 6 \times 1 & 5 \times 1 + 6 \times 0 & 5 \times 1+6\times 1 \\
    \end{array}
    \right)
    =
    \begin{pmatrix}
        2 & 1 & 3  \\
        4 & 3 & 7  \\
        6 & 5 & 11
    \end{pmatrix}
\end{equation*}



% \paragraph{転置}
Sure, I'm happy to help, but you haven't provided any Japanese text for me to translate. Could you please provide the text you want translated?
% 次に\keyterm{転置}{てんち}{transpose}と呼ばれる操作を説明します。これは計算の便宜上，よく使われる行列操作のひとつです。
Next, I will explain an operation called "transpose". This is one of the matrix operations that is often used for computational convenience.

% 大きさ$n \times m$の行列$\bm{A}$における$i$行$j$列成分を$j$行$i$列成分とする$m \times n$行列のことを，元の行列$\bm{A}$の転置とよび，$\bm{A}'$や$\bm{A}^T$と表します。行列を転ばせたようなイメージです。
The $m \times n$ matrix that swaps the $i$th row, $j$th column element in the $n \times m$ matrix $\bm{A}$ with the $j$th row, $i$th column element is called the transpose of the original matrix $\bm{A}$, and it is denoted as $\bm{A}'$ or $\bm{A}^T$. It is as if the matrix has been flipped over.

\begin{equation*}
    \bm{A}=
    \begin{pmatrix}
        a_{11} & a_{12} & \cdots & a_{1m} \\
        a_{21} & a_{22} & \cdots & a_{2m} \\
        \vdots & \vdots & \ddots & \vdots \\
        a_{n1} & a_{n2} & \cdots & a_{nm} \\
    \end{pmatrix}
    \mbox{のとき，}
    \bm{A}'=
    \begin{pmatrix}
        a_{11} & a_{21} & \cdots & a_{n1} \\
        a_{12} & a_{22} & \cdots & a_{n2} \\
        \vdots & \vdots & \ddots & \vdots \\
        a_{1m} & a_{2m} & \cdots & a_{nm} \\
    \end{pmatrix}
\end{equation*}



% ベクトルも転置でき，行ベクトルを転置すると列ベクトルに，列ベクトルを転置すると行ベクトルになります。
You can transpose a vector, when a row vector is transposed it becomes a column vector, and when a column vector is transposed it becomes a row vector.

\begin{equation*}
    \bm{a} =
    \begin{pmatrix}
        a_1 & a_2 & \cdots & a_n
    \end{pmatrix}
    \mbox{のとき，}
    \bm{a}'=
    \begin{pmatrix}
        a_1    \\
        a_2    \\
        \vdots \\
        a_n
    \end{pmatrix}
\end{equation*}


% また，転置には以下のような性質があります。これは知識として知っておくだけでよいでしょう。
Also, transposition has the following properties. It is good to know this just as knowledge.
\begin{enumerate}
	\item $(\bm{A} ')' = \bm{A}$
	\item $(\bm{A} + \bm{B} )' = \bm{A}' + \bm{B}'$
	\item $(\bm{A} \bm{B})'=\bm{B}'\bm{A}'$
	\item $(c\bm{A})'=c\bm{A}'$
\end{enumerate}


% \paragraph{逆行列}
I'm sorry, but you haven't provided any Japanese text to be translated. Could you please provide the text you need translated into English?
% 最後に\keytermL{逆行列}{ぎゃくぎょうれつ}のお話をします。逆行列は割り算のイメージです。ある行列にその逆行列をかけると単位行列になる，つまり割ると1になるような行列のことです。
Lastly, I'll talk about the inverse matrix. The inverse matrix has the image of division. If you multiply a matrix by its inverse matrix, it becomes a unit matrix, that is, it's like a matrix that becomes 1 when divided.

% 正確に表現すると，ある正方行列$\bm{A}$に対し，$\bm{A}\bm{X}=\bm{I}$となるような行列$\bm{X}$が存在するとき，これを$\bm{A}$の\keyterm{逆行列}{ぎゃくぎょうれつ}{inverse}と呼び，$\bm{A}^{-1}$で表します。正方行列でない場合に逆行列はありませんし，正方行列であっても逆行列が存在しない場合もあります。逆行列の例をみてみましょう。
To put it precisely, for a given square matrix $\bm{A}$, if there exists a matrix $\bm{X}$ such that $\bm{A}\bm{X}=\bm{I}$, this is called the \keyterm{inverse}{ぎゃくぎょうれつ}{inverse} of $\bm{A}$, and is expressed as $\bm{A}^{-1}$. Inverse matrices do not exist for non-square matrices, and there are also square matrices for which no inverse matrix exists. Let's look at an example of an inverse matrix.
% $\bm{A}=\begin{pmatrix}2&1\\5&3\end{pmatrix}$とすると，次の計算が成り立ちます。
Given $\bm{A}=\begin{pmatrix}2&1\\5&3\end{pmatrix}$, the following calculation holds true.
\begin{equation*}
	\bm{AB}=\begin{pmatrix}2&1\\5&3\end{pmatrix}\begin{pmatrix}3&-1\\-5&2\end{pmatrix}=\begin{pmatrix}1&0\\0&1\end{pmatrix}
\end{equation*}

% このとき$\bm{B}$は$\bm{A}$の逆行列，すなわち$\bm{A}^{-1}=\bm{B}$といえます。
At this time, $\bm{B}$ is the inverse matrix of $\bm{A}$, that is, $\bm{A}^{-1}=\bm{B}$.

% とくに対角行列の逆行列は，対角成分の逆数をそれぞれ対角成分とする行列になります。
The inverse matrix of a diagonal matrix is a matrix whose diagonal elements are the reciprocals of the original diagonal elements.
\begin{equation*}
	\bm{D} = \begin{pmatrix} d_1    & 0      & \cdots & 0      \\
                0      & d_2    & \cdots & 0      \\
                \vdots & \vdots & \ddots & \vdots \\
                0      & 0      & \cdots & d_n\end{pmatrix}
	\mbox{のとき，}
	\bm{D}^{-1} = \begin{pmatrix} \frac{1}{d_1} & 0             & \cdots & 0             \\
                0             & \frac{1}{d_2} & \cdots & 0             \\
                \vdots        & \vdots        & \ddots & \vdots        \\
                0             & 0             & \cdots & \frac{1}{d_n}\end{pmatrix}
\end{equation*}


% 逆行列は，行列の世界の割り算のようなものです。これで一通り四則演算の定義ができました。
The inverse matrix is like division in the world of matrices. With this, we can define all four arithmetic operations.

% \section{行列を使うと便利なこと}
\section{Convenient Aspects of Using Matrices}
\label{MatrixAndEauqtion}
% さて，ここまでで行列の計算の話をして来ましたが，どこが良いのかいまいちピンとこない，という人もいるかもしれません。そこで最後にどうしてこのような計算をするのか，何が良いのかを説明してみたいと思います。
So far, we have been discussing matrix calculations, but some of you may not yet understand why they are beneficial. Therefore, I would like to conclude by explaining why we perform these calculations and why they are advantageous.
% \subsection{行列と方程式}
\subsection{Matrices and Equations}
% 線形代数は「便利な書き方」の学問です。便利な書き方をするためにルールが作られていますから，ルールから学ぶと「なんでそんな変な操作をするんだ」という気持ちになるのもわかります。
Linear algebra is the study of "convenient notation." I understand you may feel "why am I doing such strange operations," because the rules are made up in order to make notation convenient.

% では何が便利になるのでしょうか。これは方程式を解くことと関係があります。
So what will become convenient? This has to do with solving equations.
% たとえば，以下のような連立方程式があったとしましょう。
For example, let's suppose we have a system of equations like the following.

\[
	\left \{
	\begin{array}{lcl}
		x-2y-5z  & = & 3 \\
		5x+4y+3z & = & 1 \\
		3x+y-3z  & = & 6
	\end{array}
	\right .
\]
% これは行列で表現すると，次のようになります。
When expressed as a matrix, it becomes as follows.
\[
	\begin{pmatrix} 1 & -2 & -5 \\
                5 & 4  & 3  \\
                3 & 1  & -3\end{pmatrix}
	\begin{pmatrix} x \\ y \\ z \end{pmatrix}
	=
	\begin{pmatrix} 3 \\ 1 \\ 6 \end{pmatrix}
\]

% この左辺を行列とベクトルの式の計算ルールにのっとって展開してみてください。ちゃんと最初の連立方程式の左辺になることがわかると思います。かけて足して，という面倒な計算ルールは，連立方程式を簡単に表記するためのものだったのですね。
Please try to expand this left side according to the calculation rules of matrices and vectors. I think you will see that it properly becomes the left side of the first set of equations. The troublesome calculation rules of multiplying and adding were for the purpose of simply notating the system of equations.

% 最終的にはこの方程式を解いて，次のように答えを求めます。
Ultimately, solve this equation then find the answer as follows.
\[
	\begin{pmatrix} 1 & -2 & -5 \\
                5 & 4  & 3  \\
                3 & 1  & -3\end{pmatrix}
	\begin{pmatrix} x=-1 \\ y=3 \\ z=-2 \end{pmatrix}
	=
	\begin{pmatrix} 3 \\ 1 \\ 6 \end{pmatrix}
\]


% 皆さんも学校で習ったように，このような連立方程式を解く方法として，加減法や代入法というのがあります。ですがここはひとつ，行列を使った解法を考えて見ましょう。
As everyone learned in school, there are methods such as the addition-subtraction method and the substitution method for solving such systems of equations. However, let's consider a method of solving using matrices here.

% そのような解法のひとつ，消去法は，ひとつの方程式を何倍かして，他の方程式に加えることにより，方程式をどんどん簡単にしていくというものです。
One such solution, the elimination method, is to multiply one equation by a certain amount and add it to another equation to simplify the equation gradually.
% まず，第一の式を5倍，あるいは3倍して，第二，第三の式からxの項を消去します。
First, multiply the first equation by 5 or 3, then eliminate the x term from the second and third equations.
\[ \left \{
	\begin{array}{rcl}
		x-2y-5z  & = & 3  \\
		-14y-28z & = & 14 \\
		-7y-12z  & = & 3
	\end{array}
	\right .
\]
% 第二の式の係数を簡単にしておきましょう。
Let's simplify the coefficient of the second equation.
\[ \left \{
	\begin{array}{rcl}
		x-2y-5z & = & 3  \\
		y+2z    & = & -1 \\
		-7y-12z & = & 3
	\end{array}
	\right .
\]
% 第二の式を7倍して，第三の式からyを消去します。
Multiply the second equation by 7, and eliminate y from the third equation.
\[ \left \{
	\begin{array}{rcl}
		x-2y-5z & = & 3  \\
		y+2z    & = & -1 \\
		2z      & = & -4
	\end{array}
	\right .
\]
% あとはこれの3行目から$z=-2$が得られ，芋づる式に$x=-1$，$y=3$が得られました。
From the third line of this, $z=-2$ was obtained, and $x=-1$, $y=3$ were obtained in a chain reaction.


% この操作は，式を一本ずつ，あるいは2つの式を組み合わせて文字を消していく消去法を係数全体に行う操作になっています。実際，ここで操作される係数だけ見ていくと，次のようになります。
This operation involves applying the elimination method, which eliminates characters one equation at a time or combines two equations, to all the coefficients. In fact, if we only look at the coefficients being operated on here, it will be as follows.

% \paragraph{第一段階}
Sorry, but you didn't provide any Japanese text for me to translate. Could you please provide the text you want to be translated?
\[
	\begin{pmatrix} 1 & -2 & -5  \\
                0 & 1  & 2   \\
                0 & -7 & -12\end{pmatrix}
	\begin{pmatrix} x \\ y \\ z \end{pmatrix}
	=
	\begin{pmatrix} 3 \\ -1 \\ 3 \end{pmatrix}
\]

% \paragraph{第二段階}
You didn't provide any Japanese text. Could you please provide the text you want to be translated?

\[
	\begin{pmatrix} 1 & -2 & -5 \\
                0 & 1  & 2  \\
                0 & 0  & 2\end{pmatrix}
	\begin{pmatrix} x \\ y \\ z \end{pmatrix}
	=
	\begin{pmatrix} 3 \\ -1 \\ -4 \end{pmatrix}
\]

% さらにこの方法を改良した，ガウス−ジョルダンの消去法というものがあります。この手法による係数の変化を，行列表記で見ていくことにします。
Moreover, there is a method called Gauss-Jordan elimination, which is an improved version of this method. We will look at the changes in coefficients using this method in a matrix notation.

% まず第一段目は同じです。
Firstly, the first paragraph is the same.
\[
	\begin{pmatrix} 1 & -2 & -5  \\
                0 & 1  & 2   \\
                0 & -7 & -12\end{pmatrix}
	\begin{pmatrix} x \\ y \\ z \end{pmatrix}
	=
	\begin{pmatrix} 3 \\ -1 \\ 3 \end{pmatrix}
\]
% 次に，第二の方程式を用いて第一と第三の式から$y$の係数を消してしまいます。
Next, use the second equation to eliminate the coefficient of $y$ from the first and third equations.
\[
	\begin{pmatrix} 1 & 0 & -1 \\
                0 & 1 & 2  \\
                0 & 0 & -2\end{pmatrix}
	\begin{pmatrix} x \\ y \\ z \end{pmatrix}
	=
	\begin{pmatrix} 1 \\ -1 \\ 4 \end{pmatrix}
\]
% 最後に，第三の式の$z$の係数を1にして，第一，第二式の$z$の係数を消してしまいましょう。
Finally, let's make the coefficient of $z$ in the third equation 1 and eliminate the coefficient of $z$ in the first and second equations.
\[
	\begin{pmatrix} 1 & 0 & 0 \\
                0 & 1 & 0 \\
                0 & 0 & 1\end{pmatrix}
	\begin{pmatrix} x \\ y \\ z \end{pmatrix}
	=
	\begin{pmatrix} -1 \\ 3 \\ -2 \end{pmatrix}
\]

% 最後の形を見ると，左辺は単位行列になっていますから，
Looking at the final form, the left side has become a unit matrix, so
\[
	\begin{pmatrix}
		x \\
		y \\
		z
	\end{pmatrix}
	=
	\begin{pmatrix}
		-1 \\
		3  \\
		-2
	\end{pmatrix}
\]
% と解が求められたことがわかります。
It has been understood that the solution was required.
% ここで注目すべきは，連立方程式の解を求めるプロセスは係数行列を単位行列に変えていくプロセスだった，ということです。係数行列が単位行列になれば，それはもう答えを出したことになるのです。
The key point here is that the process of finding the solution to a system of equations was a process of turning the coefficient matrix into an identity matrix. If the coefficient matrix becomes an identity matrix, it means that an answer has already been produced.

% さて，係数行列を$\bm{A}$とすると，その逆行列$\bm{A}^{-1}$があれば$\bm{A}^{-1}\bm{A} = \bm{I}$となるのでした。であれば，連立方程式の右辺にあったベクトルに$\bm{A}^{-1}$をかけてやれば，一気に答えが求まるではないですか。
Now, let's denote the coefficient matrix as $\bm{A}$. If its inverse matrix $\bm{A}^{-1}$ exists, we have $\bm{A}^{-1}\bm{A} = \bm{I}$. In that case, if we multiply the vector on the right-hand side of the system of equations by $\bm{A}^{-1}$, we can find the answer all at once, can't we?

% 実際に見て見ましょう。
Let's actually take a look.
\[
	\begin{pmatrix} 1 & -2 & -5 \\
                5 & 4  & 3  \\
                3 & 1  & -3\end{pmatrix}
	\begin{pmatrix} x \\ y \\ z \end{pmatrix}
	=
	\begin{pmatrix} 3 \\ 1 \\ 6 \end{pmatrix}
\]
% この連立方程式に対して，次のような操作をします。
We will perform the following operations on this system of equations.
\[
	\begin{pmatrix} 1 & -2 & -5 \\
                5 & 4  & 3  \\
                3 & 1  & -3\end{pmatrix}^{-1}
	\begin{pmatrix} 1 & -2 & -5 \\
                5 & 4  & 3  \\
                3 & 1  & -3\end{pmatrix}
	\begin{pmatrix} x \\ y \\ z \end{pmatrix}
	=
	\begin{pmatrix} 1 & -2 & -5 \\
                5 & 4  & 3  \\
                3 & 1  & -3\end{pmatrix}^{-1}
	\begin{pmatrix} 3 \\ 1 \\ 6 \end{pmatrix}
\]

% とします\footnote{数値的には
You didn't provide the Japanese text that you want to translate into English. Could you please provide the text?
% $\begin{pmatrix} 1 & -2 & -5 \\
I'm sorry, there seems to be a misunderstanding. As an AI model, I do not have the ability to translate mathematical equations or TeX code from Japanese to English. However, I can translate text. Please feel free to provide any Japanese text, and I'd be more than happy to assist you.
%                 5 & 4  & 3  \\
I'm sorry, but the text provided seems to be an incomplete code snippet, and not Japanese text. Please verify and provide the correct information so I can assist you better.
%                 3 & 1  & -3\end{pmatrix}^{-1}
I'm sorry but I need the Japanese text to translate. The text you provided appears to be formatted for a math equation. Please provide the correct Japanese text.
%     =
I'm sorry but you didn't provide any Japanese text to translate into English. Please provide the text you'd like me to translate.
%     \begin{pmatrix} 15/28 & 11/28 & -1/2 \\
I'm sorry for misunderstanding, but I can't translate mathematical expressions or codes. Could you please provide the Japanese text that you want translated into English, thanks.
%                 -6/7  & -3/7  & 1    \\
Sure, I would love to help. However, you've provided a matrix element written in latex format, not Japanese text. We can't translate mathematical expressions. Please provide the Japanese text you want to be translated.
%                 1/4   & 1/4   & -1/2\end{pmatrix}
I'm afraid there is a mistake as the given text isn't actually Japanese, but it appears to be a piece of coding or LaTeX. This is not a piece of text that can be translated into English.
% $
Sorry, but you didn't provide any Japanese text to be translated. Please provide the text.
% という行列です}。すると左辺は単位行列になりますから，次のように計算すれば一気に答えが求まることになるのです。
This translates to: "This is the matrix. Then, since the left side becomes the unit matrix, if you calculate as follows, you will get the answer all at once."
\[
	\begin{pmatrix} 1 & 0 & 0 \\
                0 & 1 & 0 \\
                0 & 0 & 1\end{pmatrix}
	\begin{pmatrix} x \\ y \\ z \end{pmatrix}
	=
	\begin{pmatrix} -1 \\ 3 \\ -2 \end{pmatrix}
\]
% つまり，連立方程式を解くという問題が，係数行列の逆行列を求める問題になります。また，逆行列は存在しないこともある，ということでしたが，その場合その連立方程式は解けない，ということになります。
In other words, the problem of solving a system of equations becomes the problem of finding the inverse matrix of the coefficient matrix. It was also mentioned that the inverse matrix may not exist. In that case, it means that the system of equations cannot be solved.


% \section{課題}
\section{Task}
\paragraph{線形代数の練習問題}行列$\bm{A}=\begin{pmatrix}1 & 2 & -1\\3 & 1 & 0\end{pmatrix}$のとき，次の計算をしなさい。なお，$\bm{I}_n$とは$n \times n$の単位行列，$\bm{O}$とはすべての要素が0の適当なサイズの正方行列であることを表します。
\begin{enumerate}
    \item $\bm{A}'\bm{A}$
    \item $\bm{AA}'$
    \item $\bm{A}\bm{I}_3$
    \item $\bm{A}$
\end{enumerate}
\paragraph{線形代数の練習問題その2}行列$\bm{A}=\begin{pmatrix}1 & 3 \\2 & 4\end{pmatrix},\bm{B}=\begin{pmatrix}3&1&5\\2&4&6\end{pmatrix},\bm{C}=\begin{pmatrix}4&1\\2&3\end{pmatrix}$，列ベクトル$\bm{x}=\begin{pmatrix}1\\2\\3\end{pmatrix}$,行ベクトル$\bm{y}=\begin{pmatrix}2&8\end{pmatrix}$とするとき，次の計算をしなさい。なお，計算が定義されていないものについては「計算できない」と回答しなさい。

\begin{enumerate}
	\item $\bm{A}+\bm{B}$
	\item $\bm{A}-\bm{C}$
	\item $\bm{AB}$
	\item $\bm{AC}$
	\item $\bm{B}'\bm{A}$
	\item $\bm{A}\bm{y}'$
	\item $\bm{xy}$
	\item $\bm{xB}$
	\item $\bm{x}'\bm{B}'$
	\item $\bm{yx}$
\end{enumerate}


% \paragraph{連立方程式を解く}次の連立方程式を解きなさい。
\paragraph{Solving Simultaneous Equations}Please solve the following simultaneous equations.
\[
	\left\{\begin{array}{rr}x-2 y+3 z= & 1 \\ 3 x+y-5 z= & -4 \\ -2 x+6 y-9 z= & -2\end{array}\right.
\]

\end{document}


